\documentclass[aspectratio=169,11pt]{beamer}
\usetheme{default}
\setbeamertemplate{navigation symbols}{}
\setbeamertemplate{footline}{%
  \hfill\insertframenumber/\inserttotalframenumber\hspace*{4pt}\vspace{2pt}%
}
\usepackage{lmodern}
\usepackage[most]{tcolorbox}
\usepackage{listings}
\usepackage{tikz}
\usetikzlibrary{arrows.meta,positioning,shapes.geometric,fit,backgrounds}
\usepackage{booktabs}
\usepackage{hyperref}
\usepackage{amssymb}

\definecolor{lcgreen}{HTML}{2E7D32}
\definecolor{lcblue}{HTML}{1565C0}
\definecolor{lcorange}{HTML}{EF6C00}
\definecolor{lcred}{HTML}{C62828}
\definecolor{codebg}{HTML}{F5F5F5}
\definecolor{docgray}{gray}{0.55}

\newtcolorbox{conceptbox}[1][]{colback=yellow!20,colframe=yellow!60!black,title=#1,fonttitle=\bfseries,top=2pt,bottom=2pt,left=4pt,right=4pt}
\newtcolorbox{warningbox}[1][]{colback=red!15,colframe=red!60!black,title=#1,fonttitle=\bfseries,top=2pt,bottom=2pt,left=4pt,right=4pt}
\newtcolorbox{tipbox}[1][]{colback=green!10,colframe=green!50!black,title=#1,fonttitle=\bfseries,top=2pt,bottom=2pt,left=4pt,right=4pt}
\newtcolorbox{codebox}[1][]{colback=blue!10,colframe=blue!40!black,title=#1,fonttitle=\bfseries,top=2pt,bottom=2pt,left=4pt,right=4pt}

\lstset{basicstyle=\ttfamily\tiny,language=Python}

\newcommand{\docref}[1]{{\scriptsize\color{docgray}[Docs: #1]}}

\title{Section 10: Subgraphs \& Modularity}
\subtitle{Composing Complex Systems from Smaller Graphs}
\author{LangGraph v1.0.x}
\date{March 2026}

\begin{document}

\frame{\titlepage}

\begin{frame}
\frametitle{Mental Model: Subgraphs}
\begin{conceptbox}[Core Idea]
\setlength{\itemsep}{1pt}
\begin{itemize}
\item Build complex systems from smaller graphs
\item Reuse graph definitions like functions
\item Manage complexity through composition
\item Team-friendly: teams own subgraphs
\end{itemize}
\end{conceptbox}
\vspace{0.1cm}
\begin{tipbox}[Analogy]
Like function decomposition: main orchestrates sub-tasks.
\end{tipbox}
\docref{Subgraphs Guide}
\end{frame}

\begin{frame}
\frametitle{Why Subgraphs Matter}
\begin{columns}[T]
\column{0.48\textwidth}
\begin{warningbox}[Monolithic Graph]
\setlength{\itemsep}{1pt}
\begin{itemize}
\item 50+ nodes in one file
\item Hard to understand
\item Difficult to test
\item Team scaling hard
\end{itemize}
\end{warningbox}
\column{0.48\textwidth}
\begin{tipbox}[With Subgraphs]
\setlength{\itemsep}{1pt}
\begin{itemize}
\item 3-4 nodes per graph
\item Clear responsibility
\item Easy unit tests
\item Team ownership
\end{itemize}
\end{tipbox}
\end{columns}
\end{frame}

\begin{frame}
\frametitle{Defining a Subgraph}
\begin{tipbox}[Create \& Compile]
\texttt{builder = StateGraph(SearchState)}\\
\texttt{builder.add\_node("search", search\_node)}\\
\texttt{builder.add\_edge("search", "filter")}\\
\texttt{search\_graph = builder.compile()}
\end{tipbox}
\docref{StateGraph}
\end{frame}

\begin{frame}
\frametitle{Adding Subgraph as Node}
\begin{tipbox}[Parent Graph]
\texttt{parent.add\_node("search\_task", search\_graph)}\\
\texttt{parent.add\_node("research", research\_graph)}\\
\texttt{parent.add\_edge("search\_task", "research")}
\end{tipbox}
\vspace{0.1cm}
\begin{tipbox}[Key Point]
Compiled graph can be a node in parent.
\end{tipbox}
\docref{add\_node()}
\end{frame}

\begin{frame}
\frametitle{State Mapping: Parent $\leftrightarrow$ Child}
\begin{warningbox}[Challenge]
\setlength{\itemsep}{1pt}
\begin{itemize}
\item Parent: \texttt{\{query, user\_id, results\}}
\item Child: \texttt{\{query, results\}}
\item Need to map between schemas
\end{itemize}
\end{warningbox}
\vspace{0.1cm}
\begin{tipbox}[Solution: Shared Keys]
\setlength{\itemsep}{1pt}
\begin{itemize}
\item Common keys auto-sync
\item Child uses only known fields
\item Parent sees child output changes
\end{itemize}
\end{tipbox}
\docref{State Mapping}
\end{frame}

\begin{frame}
\frametitle{Shared State Keys}
\begin{tipbox}[Parent \& Child Share]
\texttt{class ParentState(TypedDict):}\\
\texttt{query: str; results: list; user\_id: str}\\
\texttt{class ChildState(TypedDict):}\\
\texttt{query: str; results: list}
\end{tipbox}
\vspace{0.1cm}
\begin{tipbox}[Behavior]
Child reads query, modifies results. Parent sees both.
\end{tipbox}
\end{frame}

\begin{frame}
\frametitle{State Transformation}
\begin{tipbox}[Custom Mapper]
\texttt{from langgraph.graph import compose}\\
\texttt{mapping = \{"parent\_query": "child\_query"\}}\\
\texttt{child\_node = compose(child\_graph,}\\
\texttt{input\_map=mapping, output\_map=mapping)}
\end{tipbox}
\vspace{0.1cm}
\begin{tipbox}[When Needed]
Parent/child have very different schemas.
\end{tipbox}
\docref{compose()}
\end{frame}

\begin{frame}
\frametitle{Complete Example: Research Agent}
\begin{tipbox}[Web Search Subgraph]
\texttt{class SearchState(TypedDict):}\\
\texttt{query: str; sources: list}\\
\texttt{search\_builder.add\_node("search", web\_search)}\\
\texttt{search\_graph = search\_builder.compile()}
\end{tipbox}
\vspace{0.1cm}
\begin{tipbox}[Parent Graph]
\texttt{parent.add\_node("research", search\_graph)}\\
\texttt{parent.add\_node("write", write\_report)}
\end{tipbox}
\end{frame}

\begin{frame}
\frametitle{Parent Graph Diagram}
\begin{tipbox}[Hierarchy]
Entry -> Subgraph (compiled) → Exit
\end{tipbox}
\vspace{0.1cm}
\begin{conceptbox}[Key Concept]
Parent passes input to subgraph, receives output back.
\end{conceptbox}
\end{frame}

\begin{frame}
\frametitle{Nested Subgraphs}
\begin{conceptbox}[Multi-Level Composition]
\setlength{\itemsep}{1pt}
\begin{itemize}
\item Level 1: Orchestrator graph
\item Level 2: Multi-task subgraphs
\item Level 3: Leaf subgraphs
\item No hard depth limit
\end{itemize}
\end{conceptbox}
\vspace{0.1cm}
\begin{tipbox}[Design Principle]
Keep depth 2-3 levels for maintainability.
\end{tipbox}
\end{frame}

\begin{frame}
\frametitle{Checkpointing Across Subgraphs}
\begin{tipbox}[Unified Checkpointer]
\texttt{saver = SqliteSaver(":memory:")}\\
\texttt{compiled = parent.compile(checkpointer=saver)}
\end{tipbox}
\vspace{0.1cm}
\begin{tipbox}[Key Insight]
\setlength{\itemsep}{1pt}
\begin{itemize}
\item Same checkpointer for all levels
\item All state changes tracked in DB
\item Resume from any checkpoint
\end{itemize}
\end{tipbox}
\docref{Checkpointer}
\end{frame}

\begin{frame}
\frametitle{Streaming from Subgraphs}
\begin{tipbox}[Parent Streams Child Events]
\texttt{for event in parent.stream(input\_data):}\\
\texttt{print(f"Parent event: \{event\}")}
\end{tipbox}
\vspace{0.1cm}
\begin{tipbox}[Behavior]
\setlength{\itemsep}{1pt}
\begin{itemize}
\item Events from subgraph bubble up
\item Parent sees all node outputs
\item Can filter by node name
\end{itemize}
\end{tipbox}
\end{frame}

\begin{frame}
\frametitle{Human-in-the-Loop in Subgraphs}
\begin{tipbox}[Interrupt Node]
\texttt{user\_approval = get\_input(...)}\\
\texttt{if not user\_approval:}\\
\texttt{raise GraphInterrupt(...)}
\end{tipbox}
\vspace{0.1cm}
\begin{tipbox}[Integration]
\setlength{\itemsep}{1pt}
\begin{itemize}
\item Interrupt works in subgraph nodes
\item Parent resumes from checkpoint
\item Same thread\_id across levels
\end{itemize}
\end{tipbox}
\docref{GraphInterrupt}
\end{frame}

\begin{frame}
\frametitle{Modular Agent Architecture Pattern}
\begin{conceptbox}[Typical Structure]
\setlength{\itemsep}{1pt}
\begin{itemize}
\item Orchestrator: routes to subagents
\item Search Agent: web search subgraph
\item Writer Agent: content synthesis
\item Reviewer Agent: fact-checking
\end{itemize}
\end{conceptbox}
\vspace{0.1cm}
\begin{tipbox}[Benefits]
\setlength{\itemsep}{1pt}
\begin{itemize}
\item Each team owns subgraph
\item Parallel development
\item Easy test and deploy
\end{itemize}
\end{tipbox}
\docref{Agent Patterns}
\end{frame}

\begin{frame}
\frametitle{Common Subgraph Pitfalls}
\begin{warningbox}[State Key Conflicts]
Parent/child expect different types for same key.\\
\textbf{Fix:} Define clear contracts with TypedDict.
\end{warningbox}
\vspace{0.08cm}
\begin{warningbox}[Checkpoint Isolation]
Each subgraph uses own checkpointer.\\
\textbf{Fix:} Pass single checkpointer to compile().
\end{warningbox}
\vspace{0.08cm}
\begin{warningbox}[Tight Coupling]
Subgraph depends on parent structure.\\
\textbf{Fix:} Use input\_map, output\_map for mapping.
\end{warningbox}
\end{frame}

\begin{frame}
\frametitle{Summary: Subgraphs Cheat Sheet}
\begin{conceptbox}[Key Points]
\setlength{\itemsep}{1pt}
\begin{itemize}
\item Define: \texttt{builder.compile()}
\item Add to parent: \texttt{add\_node("name", subgraph)}
\item Shared keys auto-sync state
\item \texttt{compose()} for complex mapping
\item Same checkpointer for all children
\item Streaming bubbles up
\item Interrupt works in subgraph nodes
\item Keep 2-3 levels deep
\end{itemize}
\end{conceptbox}
\end{frame}

\begin{frame}
\frametitle{Self-Check Questions}
\begin{enumerate}
\setlength{\itemsep}{2pt}
\item Why use subgraphs vs one large?
\item How add compiled graph as node?
\item Parent vs child state difference?
\item How shared keys work?
\item When use \texttt{compose()}?
\item Checkpoints across boundaries?
\item Streaming bubbles up?
\item Nesting depth recommendation?
\end{enumerate}
\end{frame}

\end{document}
